%% ============================================================
\section{AI Integration and LLM Pipeline}
\label{sec:ai}
%% ============================================================

\subsection{Design Goals}

The AI subsystem of \echoray{} is designed to satisfy three requirements:

\begin{enumerate}[leftmargin=*, label=\textbf{R\arabic*.}]
  \item \textbf{Low Latency.} Code-generation responses must be delivered
    within a perceptually acceptable window ($<$3\,s) to maintain the
    conversational flow of a collaborative session.
  \item \textbf{Structured Output.} AI responses must be machine-parseable
    JSON so that the front end can apply syntax highlighting, diff views, and
    code insertion without unreliable regex extraction.
  \item \textbf{Broadcast Consistency.} All project members must receive the
    identical AI response simultaneously, ensuring shared context.
\end{enumerate}

\subsection{AI Service Implementation}

The AI service is encapsulated in \texttt{server/src/services/ai.service.ts}.

\begin{lstlisting}[caption={AI service---Gemini configuration and query method
  (\texttt{ai.service.ts}).}, label=lst:aiservice]
import { GoogleGenerativeAI } from "@google/generative-ai";

const genAI = new GoogleGenerativeAI(process.env.AI_API_KEY!);

const model = genAI.getGenerativeModel({
  model: "gemini-2.5-flash",
  generationConfig: {
    responseMimeType: "application/json",
    temperature: 0.4,
  },
  systemInstruction: `You are an expert JavaScript developer with
    10+ years of experience in MERN stack. When asked to provide code
    or explanation, always respond in JSON format:
    { "text": "<explanation>", "fileTree": { "<filename>": {
      "file": { "contents": "<code>" } }, ... } }`,
});

export const generateResult = async (prompt: string): Promise<string> => {
  const result = await model.generateContent(prompt);
  return result.response.text();
};
\end{lstlisting}

Several implementation choices warrant discussion:

\textbf{Model selection (gemini-2.5-flash).}
Gemini~2.5~Flash is Google's latency-optimised variant of the Gemini~2.5
family, providing a favourable latency--quality trade-off for interactive
applications. The model supports JSON-mode output via
\texttt{responseMimeType: "application/json"}, eliminating prompt-engineering
gymnastics required to coerce older models into structured output.

\textbf{Temperature ($T = 0.4$).}
A temperature below the default ($T = 1.0$) biases the model towards more
deterministic, syntactically predictable output. This is appropriate for code
generation, where correctness and reproducibility are more important than
creative diversity. Wei \etal~\cite{wei2022chain} and subsequent calibration
studies corroborate that $T \in [0.2, 0.6]$ yields the best code-generation
accuracy.

\textbf{System instruction.}
The system instruction anchors the model to a specific persona (experienced
JavaScript/MERN developer) and enforces the output schema. The schema
stipulates two fields: \texttt{text} (a human-readable explanation) and
\texttt{fileTree} (a file tree matching Equation~\ref{eq:filetree}). This
design means the AI response is directly mountable into the WebContainer
instance, enabling one-click execution of AI-generated code.

\subsection{AI Controller and Route}

The REST endpoint (\texttt{GET /api/v1/ai/get-result?prompt=...}) is handled
by \texttt{ai.controller.ts}:

\begin{lstlisting}[caption={AI controller (\texttt{ai.controller.ts}).},
  label=lst:aictl]
export const getResult = async (req: Request, res: Response) => {
  const prompt = req.query.prompt as string;
  const result = await generateResult(prompt);
  res.json({ result });
};
\end{lstlisting}

This endpoint is intentionally unauthenticated to allow direct programmatic
access; the primary invocation path flows through Socket.io where project
membership is already verified.

\subsection{AI-Broadcast Flow}

\Cref{fig:ai_flow} illustrates the complete end-to-end flow from a user
typing \texttt{@ai generate a REST endpoint for user login} to all
collaborators receiving the generated code.

\begin{figure}[H]
  \centering
  \begin{tikzpicture}[
    font=\small,
    node distance=0.45cm and 0.1cm,
    step/.style={
      rectangle, rounded corners=3pt,
      minimum width=3cm, minimum height=0.7cm,
      text centered, draw, thick, fill=white
    },
    timeline/.style={-Stealth, thick, gray},
    arr/.style={-Stealth, thick, blue!70!black},
  ]
    %% ── Swim-lane headers ────────────────────────────────────
    \node[font=\bfseries\small] (hdr_u)  at (-4.5, 0.4) {User A};
    \node[font=\bfseries\small] (hdr_s)  at (0,    0.4) {Server};
    \node[font=\bfseries\small] (hdr_g)  at (4.5,  0.4) {Gemini};
    \node[font=\bfseries\small] (hdr_b)  at (-4.5,-7.6) {All Members};

    %% ── Timeline bars ────────────────────────────────────────
    \draw[gray!40, very thick] (-4.5,0) -- (-4.5,-7.2);
    \draw[gray!40, very thick] (0,0)    -- (0,-7.2);
    \draw[gray!40, very thick] (4.5,0)  -- (4.5,-7.2);

    %% ── Steps ────────────────────────────────────────────────
    \node[step, fill=blue!10]   (t1) at (-4.5,-0.8) {Type \texttt{@ai ...}};
    \node[step, fill=blue!10]   (t2) at (-4.5,-1.8) {emit("project-message")};
    \node[step, fill=green!10]  (t3) at (0,   -2.5) {Detect \texttt{@ai} prefix};
    \node[step, fill=green!10]  (t4) at (0,   -3.4) {Extract prompt};
    \node[step, fill=green!10]  (t5) at (0,   -4.2) {Broadcast raw msg};
    \node[step, fill=orange!10] (t6) at (4.5, -4.9) {Process \& generate};
    \node[step, fill=orange!10] (t7) at (4.5, -5.8) {Return JSON response};
    \node[step, fill=green!10]  (t8) at (0,   -6.2) {\texttt{io.to(room).emit}};
    \node[step, fill=blue!10]   (t9) at (-4.5,-7.0) {Render AI message};

    %% ── Arrows ───────────────────────────────────────────────
    \draw[arr] (t2.east)  -- (t3.west)  node[midway,above]{\tiny WSS};
    \draw[arr] (t3.south) -- (t4.north);
    \draw[arr] (t4.east)  -- (t6.west)  node[midway,above]{\tiny HTTPS};
    \draw[arr] (t4.south) -- (t5.north);
    \draw[arr] (t7.west)  -- (t8.east)  node[midway,above]{\tiny JSON};
    \draw[arr] (t8.west)  -- (t9.east)  node[midway,above]{\tiny WSS};
  \end{tikzpicture}
  \caption{Sequence diagram of an \texttt{@ai} invocation. The message is
    first rebroadcast to the room (so peers see the user's query), then the
    Gemini response is broadcast to all room members including the invoker.}
  \label{fig:ai_flow}
\end{figure}

\subsection{Front-End AI Response Rendering}

The front-end \texttt{ChatBox} component (\texttt{components/ChatBox.tsx})
renders chat messages using \texttt{react-markdown} with the
\texttt{remark-gfm} plugin, providing GitHub-Flavoured Markdown support
including fenced code blocks with syntax highlighting. AI responses are
identified by \texttt{sender.email === "AI"} and styled with a distinct badge,
enabling users to visually distinguish AI-generated content from peer messages.

\subsection{Security Considerations for AI Endpoints}

\textbf{Prompt injection.}
The system instruction partially mitigates prompt injection by constraining
the model's response schema; arbitrary instruction overrides are unlikely to
produce non-JSON output that bypasses downstream parsing. Production hardening
should apply a dedicated prompt-injection detection layer~\cite{greshake2023indirect}.

\textbf{Cost management.}
The Gemini API is billed per token. Without rate limiting, a malicious actor
could exhaust the API quota. The current implementation lacks per-user rate
limiting; \cref{sec:conclusion} lists this as a priority item.

\textbf{Data privacy.}
User prompts are transmitted to Google's infrastructure. \echoray{}'s
deployment documentation should advise operators to review Google's data
processing terms and consider on-premises alternatives (e.g., Ollama
\cite{ollama2024}) for sensitive environments.
